\section{结论}\label{sec:conclusion}

本文针对多源融合机器人定位及任务优化问题，建立了四层递进的数学模型体系，主要结论如下：

\textbf{（1）问题一}：采用三次样条插值 + 一维黄金分割搜索，在 $[-250, -150]$\,s 范围内精确估计时间偏差 $\Delta\tau^* = \boldsymbol{-198.43}$\,s，对齐后 RMSE = $\boldsymbol{8.39\times10^{-11}}$\,m（达浮点精度极限），验证了无噪声条件下方法的精确性。输出 10\,Hz 轨迹 7004 点。

\textbf{（2）问题二}：建立联合最小二乘模型同时估计三参数，得 $\Delta\tau^* = \boldsymbol{-48.44}$\,s、$\Delta x^* = \boldsymbol{-3.44}$\,m、$\Delta y^* = \boldsymbol{1.81}$\,m，RMSE = $\boldsymbol{1.56}$\,m。常加速度卡尔曼滤波异步融合输出 8513 点 10\,Hz 轨迹。

\textbf{（3）问题三}：Wald 假设检验 $W=2.07$，$p=\boldsymbol{0.356} > 0.05$；AIC 对比 $\Delta\text{AIC} = \boldsymbol{-2.81} < 2$。双判据一致判定\textbf{无显著系统偏差}（接受 $H_0$），按零偏模型融合输出 3691 点轨迹。

\textbf{（4）问题四}：SG 平滑（窗口 151 点）+ 贪心调度，在距离/速度/加速度/准备时长/角差约束下，1.2\,s 完成规划：\textbf{射击 8 次 + 拍照 16 次 = 24 个合法任务}，全部通过独立核验。

\subsection{模型推广}

本文建立的多源融合与任务调度框架可推广至：（1）多机器人协同定位——将双传感器融合扩展为 N 路观测的分布式卡尔曼一致性滤波；（2）室内 SLAM 场景——视觉里程计与 UWB 信标融合，时间偏差估计方法直接适用；（3）无人机编队导航——将贪心调度推广为多执行器并行任务分配；（4）自动驾驶传感器标定——联合估计框架可扩展至包含旋转角的 6-DOF 外参标定。
